\chapter{Workplan}
This section outlines the proposed work plan for the project, detailing the key phases and activities required to achieve the stated objectives. The plan follows a structured approach, progressing from foundational research through system design, prototype development, testing, and final documentation.

\begin{itemize}

\item \textbf{Research} \\
A literature review will be conducted to develop an understanding of gesture control, its underlying principles, and its various applications. This review will help clarify the problem space, enabling a well-defined problem statement. It will also provide a basis for comparing existing gesture recognition techniques and implementations.

\item \textbf{Concept Definition}
\begin{itemize}
    \item \textbf{Aim and Objectives} \\
    This section will clearly define the overall aim of the project, along with the specific objectives required to achieve it.

    \item \textbf{Stakeholders’ Requirements} \\
    This will identify the needs and expectations of the end users. It will define who the system is intended for and outline the key functionalities and performance expectations from a user perspective.

    \item \textbf{Engineering Requirements} \\
    The technical requirements of the system will be established here, guided by the stakeholders’ needs. These requirements will translate user expectations into measurable engineering specifications.

    \item \textbf{Drone Selection} \\
    A comparison of potential drones will be conducted, highlighting the advantages and limitations of each option. This will support the selection of the most suitable platform for the project.
\end{itemize}

\item \textbf{System Architecture Design}
\begin{itemize}
    \item \textbf{Model Selection} \\
    A suitable gesture recognition and classification approach will be selected. This includes evaluating existing methods and identifying opportunities for improvement.

    \item \textbf{Model Development} \\
    The chosen model will be developed, including training procedures and parameter tuning.

    \item \textbf{Model Optimisation} \\
    The model will be refined to improve performance, accuracy, and computational efficiency.
\end{itemize}

\item \textbf{Prototype}
\begin{itemize}
    \item \textbf{Hardware Selection} \\
    Appropriate hardware components, such as cameras and sensors, will be selected. The drone selection is treated separately, as it significantly influences other hardware choices.

    \item \textbf{Assembly} \\
    The system components will be assembled into a functional prototype, including the final test drone setup.

    \item \textbf{Hardware Testing} \\
    The hardware will be tested to ensure reliable operation and to identify and resolve any faults.

    \item \textbf{Data Collection Protocol} \\
    A structured method will be designed for safely collecting and storing data from both the drone and the gesture recognition system.

    \item \textbf{Software Implementation} \\
    Once the hardware is operational and data collection methods are in place, the software will be integrated with the hardware, and initial system testing will begin.

    \item \textbf{Software Evaluation} \\
    System performance will be evaluated, and improvements will be made based on the results.

    \item \textbf{System Implementation} \\
    The fully integrated system will be tested in various environments and across different use-case scenarios to assess robustness and adaptability.

    \item \textbf{Ethics Approval} \\
    Ethical approval will be obtained, as the system testing involves human participants.

    \item \textbf{User Testing} \\
    Participants will be asked to perform a range of gestures to evaluate how effectively the system responds to different users and gesture variations.

    \item \textbf{Results Analysis} \\
    The collected data will be analysed to assess system performance and identify areas for improvement.
\end{itemize}

\item \textbf{Documentation} \\
The project report will be continuously developed throughout the project. This ensures that all findings, results, and insights are accurately recorded and preserved.

\end{itemize}
